\documentclass[11pt]{article}
\usepackage[utf8]{inputenc}
\usepackage[T1]{fontenc}
\usepackage{amsmath,amssymb,amsthm}
\usepackage{geometry}
\geometry{margin=1in}
\usepackage{array}
\usepackage{hyperref}
\hypersetup{colorlinks=true,linkcolor=blue,urlcolor=blue,citecolor=blue}

\newtheorem{lem}{Lemma}[section]
\newtheorem{prop}[lem]{Proposition}
\theoremstyle{remark}
\newtheorem*{rem}{Remark}

\newcommand{\T}{\mathsf{T}}
\newcommand{\F}{\mathsf{F}}
\newcommand{\Zmark}{\mathsf{Z}}
\newcommand{\Nph}{\mathsf{N}}

\title{\textbf{A Paradoxical, Self-Referential System of Logics}\\[4pt]
\large Four non-classical logics as resolutions of one universal diagonal,\\
with two machine-checked poles}
\author{Vitaly Reznik}
\date{2026 --- Version 1.0.0\\[2pt]
\small Lean 4 witnesses: \texttt{github.com/inventor1975/ZTL}
(\texttt{lean/ZTL.lean}, \texttt{lean/QuantumWitness.lean})\\
\small Verify from zero: \S\ref{sec:mirror} --- two files, two commands, no dependencies}

\begin{document}
\maketitle

\begin{abstract}
We propose a reading of four non-classical logics --- classical,
intuitionistic / type-theoretic, quantum (orthomodular), and Zero-Trust
Logic (ZTL) --- as four disciplined responses to a single self-referential
construction: the universal diagonal underlying the liar, Russell, Cantor,
G\"odel and Tarski (Lawvere 1969; Yanofsky 2003). Each retains all classical
structural principles \emph{save one}, and each thereby houses the diagonal
differently. Two features are invariant across all four: non-contradiction
with its associated reductio (a retained \emph{principle}), and the diagonal
itself (a retained \emph{obstruction}). We isolate the one part of this
picture that is formal rather than expository --- the complementary sacrifice
made by ZTL and by orthomodular quantum logic --- and record it as two
propositions machine-checked in Lean~4 on the \emph{empty axiom list} (no
classical choice, no quotient, no propositional extensionality). We are
explicit throughout about what is established (all four logics; the
Lawvere--Yanofsky unification), what is the author's (ZTL; the arrangement;
the two witnesses), and what is analogy carrying no proof (the ``cycle''
itself). This is a synthesis and a reading, not a theorem, a merger, a new
foundation, or a new field.

\smallskip
\noindent\textbf{Keywords.} self-reference; Lawvere fixed point; diagonal
argument; quantum logic; orthomodular lattice; three-valued logic;
paracomplete logic; Zero-Trust Logic; universal logic; Lean~4;
machine-checked proof.

\smallskip
\noindent\textbf{MSC 2020.} 03B50; 03B60; 03G12; 18C50; 03B35.
\end{abstract}

\section{Scope}\label{sec:scope}

The claim of this note is deliberately narrow. It is \emph{not} that a new
logic or a new foundation has been discovered, nor that the four logics below
are amalgamated into one calculus --- they are provably incompatible in the
precise sense of \S\ref{sec:noamalg}. The unification of the classical
antinomies under a single fixed-point schema is due to Lawvere~\cite{lawvere}
and, in the concrete set-theoretic form used here, to
Yanofsky~\cite{yanofsky}; quantum logic is Birkhoff--von
Neumann~\cite{bvn}; the intuitionistic corner is standard; the ``general
theory of logical systems'' already exists as a field (universal logic,
abstract model theory~\cite{beziau}). The single component original to the
author is \textbf{ZTL}~\cite{ztl}, which appears here only as one corner.
What this note adds is (i) the \emph{arrangement} of the four as four
resolutions of one diagonal, and (ii) two machine-checked propositions
(\S\ref{sec:mirror}) exhibiting the sharpest pair of them as complementary.
Everything expository is marked as such; everything proved is in
\S\ref{sec:mirror} and reproducible from zero.

\section{The universal diagonal}\label{sec:diag}

The classical antinomies are one construction in several presentations.
Lawvere's fixed-point theorem makes this exact.

\begin{lem}[Lawvere~\cite{lawvere}]
In a cartesian closed category, if there is a point-surjective morphism
$\varphi : A \to Y^{A}$, then every endomorphism $g : Y \to Y$ has a fixed
point $y : 1 \to Y$ with $g \circ y = y$.
\end{lem}

Contrapositively: a \emph{fixed-point-free} endomorphism of $Y$ obstructs any
point-surjection $A \to Y^{A}$. Instantiating $Y$ and $g$:

\begin{itemize}
\item $Y = 2$, $g = \neg$ (no fixed point): no surjection $A \to 2^{A}$ ---
\textbf{Cantor};
\item the same schema, read in an internal truth object with a
fixed-point-free negation, yields the \textbf{liar} and \textbf{Russell};
\item read in a provability / representation setting it yields
\textbf{G\"odel} and \textbf{Tarski}'s undefinability
(Yanofsky~\cite{yanofsky}, \S\S3--5).
\end{itemize}

The common core is \emph{self-application against a fixed-point-free map}: an
object that names its own maps, composed with a $g$ that has no fixed point.
Following Yanofsky, we call this schema \textbf{the diagonal}. It is
\emph{pre-classical}: it is not built by classical logic but is a construction
that any logic expressive enough to talk about its own sentences must
confront. A logic that ignores it does not escape it; it trivializes on it.

\section{Four resolutions}\label{sec:four}

A logic admitting self-reference must do something disciplined with the
diagonal or collapse. There are (at least) four such disciplines in current
use, and each is \emph{classical logic minus exactly one structural
principle}:

\begin{center}
\renewcommand{\arraystretch}{1.3}
\begin{tabular}{@{}>{\raggedright}p{2.6cm}>{\raggedright}p{3.4cm}p{6.4cm}@{}}
\hline
\textbf{Logic} & \textbf{Principle dropped} & \textbf{Resolution of the diagonal} \\
\hline
Classical & none & explodes on the raw liar, then localizes truth one level up (Tarski hierarchy) \\
Intuitionistic / type-theoretic & excluded middle; unrestricted $\mathrm{Type}:\mathrm{Type}$ & bans self-typing --- the universe hierarchy; Girard's paradox is the diagonal of na\"ive type theory (Coquand~\cite{coquand}; Hurkens~\cite{hurkens}) \\
Quantum (orthomodular) & distributivity & the self-negating join lives as a genuine element; complementation stays total \\
ZTL & identity $p\to p$, excluded middle, double negation & pointwise quarantine: the liar has no fixed value; it is flagged, not exploded~\cite{ztl} \\
\hline
\end{tabular}
\end{center}

Two honesty points. First, only the ZTL row is the author's; the other three
are the standard content of their fields, curated here, not discovered.
Second, the classical row is the informative one: classical logic retains
\emph{every} structural principle, which is exactly why it is the system that
must \emph{explode} on the unlocalized liar and repair itself from the
metalanguage. Retaining everything is paid for in the hierarchy.

\section{Two invariants}\label{sec:inv}

Across the four, two things do not move --- a dual pair, one positive and one
negative.

\paragraph{The floor (a retained principle).} Non-contradiction and the
reductio it licenses, $(p \to \neg p) \to \neg p$, survive in every corner. In
the quantum lattice $x \wedge x^{\perp} = 0$; in the ZTL matrix
$\neg(p \wedge \neg p)$ is designated for every value; both are among the
propositions of \S\ref{sec:mirror}.

\paragraph{The ceiling (a retained obstruction).} The diagonal survives in
every corner, precisely because each logic is expressive enough to state it.
What differs is only the \emph{manner of housing} it (\S\ref{sec:four}).

The two are linked: \textbf{non-contradiction is the guard against the
diagonal.} A logic that surrendered non-contradiction would not tame the
diagonal but be trivialized by it (\emph{ex contradictione quodlibet}). Where
the obstruction persists, the floor stands; that is why both are invariant.

\section{No consistent amalgamation}\label{sec:noamalg}

Why an irreducible plurality rather than a hierarchy or a common refinement?

Because the retained principles are \emph{pairwise incompatible}. Quantum
logic retains $x \vee x^{\perp} = 1$; ZTL withholds $p \vee \neg p$ at the
mark --- no single valuation structure validates both a total complementation
law and its pointwise failure. Type theory demands strong normalization
(every term terminates); the ZTL solver houses non-termination as the
transient phase $\Nph$ (\S\ref{sec:mirror}). Any system retaining \emph{all
four} families of principles at once would validate distributivity \emph{and}
its negation, hence be trivial --- that system is classical logic, which pays
for retaining everything with the raw diagonal it must then hierarchically
contain. There is therefore no consistent common refinement; the four are
genuinely incomparable.

We describe this informally as an \emph{orbit} around classical logic --- the
unreachable centre that keeps all principles at the price of the liar. We
stress that ``orbit'', ``centre'' and ``cycle'' are metaphor: the content is
the pairwise incompatibility above and the concrete complementarity of
\S\ref{sec:mirror}. No claim of a formal cyclic structure is made.

\section{The ZTL--quantum mirror, machine-checked}\label{sec:mirror}

The one formalizable part of the picture is the \emph{complementary sacrifice}
made by ZTL and by orthomodular logic. Both are recorded below as
propositions verified in Lean~4 over finite carriers by \texttt{decide}, on
the \textbf{empty axiom list}, in bare Lean (no \texttt{mathlib}, no imports).

\paragraph{Pole A (ZTL).} Let $M = \{\T, \F, \Zmark\}$, with connectives
generated by the \emph{zero-trust lift}: for a Boolean $f$,
$\mathrm{lift}(f)(\vec{x}) = \T$ iff $f$ returns \texttt{true} under every
classical reading of the arguments ($\Zmark$ read as both \texttt{true} and
\texttt{false}), else $\F$. Then in $M$:

\begin{prop}[\texttt{lean/ZTL.lean}, empty axiom list]\label{propA}
\leavevmode
\begin{enumerate}
\item[(i)] \emph{non-contradiction} --- $\neg(p \wedge \neg p) = \T$ for all $p$;
\item[(ii)] \emph{distributivity in both directions} ---
$p \wedge (q \vee r) = (p\wedge q) \vee (p\wedge r)$ and dually, for all $p,q,r$;
\item[(iii)] \emph{failure of identity, excluded middle and double negation} ---
$p \to p$, $\;p \vee \neg p$, $\;\neg\neg p \to p$ are not designated for all
$p$ (each is falsified at $p = \Zmark$).
\end{enumerate}
\end{prop}

Here the third symbol $\Zmark$ is a \textbf{status mark, not a truth value}:
verdicts are two-valued, $\Zmark$ marking an unverified input (greedy
register) or, reused positionally in the solver, the transient phase
$\Nph = \text{``not yet''}$ (Kleene's undefined), which never labels an input.
ZTL is thus \textbf{paracomplete} (truth-value gaps, no gluts), not
paraconsistent~\cite{ztl}.

\paragraph{Pole B (orthomodular).} Let $\mathrm{MO}_2$ be the smallest
non-distributive orthomodular lattice: $\{0, 1\}$ together with two
complementary pairs $(a, a^{\perp}), (b, b^{\perp})$, all four proper elements
mutually incomparable, with $x \wedge x^{\perp} = 0$, $x \vee x^{\perp} = 1$.
Then:

\begin{prop}[\texttt{lean/QuantumWitness.lean}, empty axiom list]\label{propB}
\leavevmode
\begin{enumerate}
\item[(i)] \emph{non-contradiction} --- $x \wedge x^{\perp} = 0$ for all $x$;
\item[(ii)] \emph{excluded middle and double negation} ---
$x \vee x^{\perp} = 1$ and $x^{\perp\perp} = x$ for all $x$;
\item[(iii)] \emph{failure of distributivity} ---
$a \wedge (b \vee b^{\perp}) = a \wedge 1 = a$, whereas
$(a\wedge b) \vee (a\wedge b^{\perp}) = 0 \vee 0 = 0$.
\end{enumerate}
\end{prop}

\paragraph{The mirror.} ZTL retains distributivity and surrenders excluded
middle, double negation and identity; orthomodular logic retains excluded
middle and double negation and surrenders distributivity; both retain
non-contradiction. Informally: ZTL fails where a proposition meets
\emph{itself}; quantum logic fails where propositions \emph{combine}.

\begin{rem}[the honest mark]
The two ``retentions of excluded middle'' are not instances of one theorem in
one system. In $\mathrm{MO}_2$, $x \vee x^{\perp} = 1$ is a \emph{lattice
identity} and $x^{\perp\perp}=x$ holds because the orthocomplement is
involutive \emph{by the choice of an ortholattice} --- mild, structural. In
$M$, the \emph{failure} of $p \vee \neg p$ is a computation in a \emph{logical
matrix}, where validity means designation. These are facts in two different
formalisms that rhyme; the ``mirror'' is an analogy made precise on each side,
not a single cross-formalism theorem. This is the boundary between what is
proved (each proposition) and what is read (their pairing).
\end{rem}

\paragraph{Reproduction} (Lean 4.29.1, measured 2026-07-18):

\begin{verbatim}
$ lean lean/ZTL.lean             # 13 objects, no axioms,  ~1.1 s
$ lean lean/QuantumWitness.lean  #  5 objects, no axioms,  ~0.5 s
\end{verbatim}

Every \texttt{\#print axioms} line reports \emph{``does not depend on any
axioms''}, definitions included: no \texttt{Classical.choice}, no
\texttt{Quot.sound}, no \texttt{propext}. The intuitionistic and classical
corners are, by contrast, \emph{metatheorems} (underivability of excluded
middle; failure of strong normalization for $\mathrm{Type}:\mathrm{Type}$;
``classical as centre'') and are not formalized here by design; the diagonal
itself, ranging over four logics that are not objects of one theory, is
likewise not a single formal object.

\section{Provenance}\label{sec:prov}

\begin{center}
\renewcommand{\arraystretch}{1.3}
\begin{tabular}{@{}>{\raggedright}p{2.4cm}p{7.7cm}>{\raggedright\arraybackslash}p{2.3cm}@{}}
\hline
 & \textbf{Content} & \textbf{Status} \\
\hline
Established & Lawvere's fixed-point theorem and the diagonal unification (Lawvere~\cite{lawvere}; Yanofsky~\cite{yanofsky}); quantum logic and non-distributivity (Birkhoff--von Neumann~\cite{bvn}); intuitionism, the universe hierarchy, Girard's paradox (Martin-L\"of~\cite{ml}; Coquand~\cite{coquand}; Hurkens~\cite{hurkens}); universal logic as a field (B\'eziau~\cite{beziau}) & literature \\
Author's & ZTL, the zero-trust paracomplete logic~\cite{ztl}; the arrangement of the four as resolutions of one diagonal; the two Lean witnesses on the empty axiom list & synthesis + machine-checked \\
Analogy (no proof) & the ``cycle / orbit / centre''; ``the floor guards the ceiling''; ``each logic pays one principle so none pays the whole diagonal'' & structural metaphor, marked \\
\hline
\end{tabular}
\end{center}

The defensible claim is narrow: four established non-classical logics align as
four disciplined resolutions of one pre-classical diagonal, and the sharpest
pair among them --- ZTL and orthomodular logic --- make complementary,
machine-checkable sacrifices. The organizing picture is offered as a reading;
its two poles are offered as proof.

\begin{thebibliography}{99}
\bibitem{bvn} G. Birkhoff, J. von Neumann. \emph{The logic of quantum
mechanics.} Ann. of Math. \textbf{37} (1936), 823--843.
\bibitem{bochvar} D. A. Bochvar. \emph{On a three-valued logical calculus and
its application to the analysis of the paradoxes of the classical extended
functional calculus.} Mat. Sbornik \textbf{4} (1938); Eng. transl. Hist.
Philos. Logic \textbf{2} (1981), 87--112.
\bibitem{coquand} T. Coquand. \emph{An analysis of Girard's paradox.} LICS
1986, 227--236.
\bibitem{hurkens} A. G. Hurkens. \emph{A simplification of Girard's paradox.}
TLCA 1995, LNCS 902, 266--278.
\bibitem{kleene} S. C. Kleene. \emph{Introduction to Metamathematics.}
North-Holland, 1952.
\bibitem{kripke} S. Kripke. \emph{Outline of a theory of truth.} J. Philos.
\textbf{72} (1975), 690--716.
\bibitem{lawvere} F. W. Lawvere. \emph{Diagonal arguments and cartesian closed
categories.} LNM 92, Springer, 1969, 134--145.
\bibitem{lukasiewicz} J. \L ukasiewicz. \emph{On three-valued logic} (1920).
In: Selected Works, North-Holland, 1970.
\bibitem{ml} P. Martin-L\"of. \emph{Intuitionistic Type Theory.} Bibliopolis,
1984.
\bibitem{lean4} L. de Moura, S. Ullrich. \emph{The Lean 4 theorem prover and
programming language.} CADE 2021, LNCS 12699, 625--635.
\bibitem{yanofsky} N. S. Yanofsky. \emph{A universal approach to
self-referential paradoxes, incompleteness and fixed points.} Bull. Symbolic
Logic \textbf{9} (2003), 362--386.
\bibitem{ztl} V. Reznik. \emph{ZTL --- Zero-Trust Logic.} Preprint, DOI
\texttt{10.5281/zenodo.21318981} (concept).
\bibitem{beziau} J.-Y. B\'eziau (ed.). \emph{Universal Logic: An Anthology.}
Birkh\"auser, 2012.
\end{thebibliography}

\vspace{1em}
\noindent\emph{AI disclosure.} This note was written in dialogue with the AI
system Claude (Anthropic); all design decisions, the framing and final
responsibility are the human author's. The reliability of \S\ref{sec:mirror}
depends on neither the author nor the AI: both Lean files verify against the
Lean~4 kernel on the empty axiom list. Everything outside \S\ref{sec:mirror}
is offered as a reading and marked as such.

\smallskip
\noindent\emph{License.} Text CC BY 4.0; the accompanying Lean files are
MIT-licensed (repository \texttt{ZTL}).

\end{document}
